\section{Revocation exposure and acknowledgment}
\label{sec:v2revocation}
Revocation is defined per required issuer. Let $r$ be the instant the broker
commits that issuer's withdrawal, and let $P_r$ contain the issued permits that
could still cause an additional effect under its old authority. Each such
permit has a retained reservation $q_i$. Already completed effects are excluded
from additional exposure; uncertain outcomes may be included conservatively.

\subsection{Admission revocation}
Admission revocation advances the issuer epoch, marks it inactive, and removes
its approval bits in one broker transaction. Every later issuance requiring
that issuer fails until fresh live approvals exist at the new epoch. Other
issuers' independent grants continue to have their specified meaning.
Previously issued permits may remain in transit, accepted at the sink, or
uncertain after a lost reply. Withdrawal alone supplies no information about
their eventual effects.

For each $i\in P_r$, let $e_i$ be the charge of additional effects it causes
after $r$. Under the supported sink contract, identity deduplication and maximum
charge binding imply $0\le e_i\le q_i$. New old-epoch permits are not issued.
Therefore
\begin{equation}
       \sum_{i\in P_r} e_i\ \le\ \sum_{i\in P_r}q_i .
\label{eq:v2exposure}
\end{equation}
The core includes the corresponding finite-sequence subset-sum lemma. This
bound counts charged effects, not wall-clock duration, number of model tokens,
or the monetary value of every possible external consequence. Its validity
requires that the normalized maximum charge really bounds the sink effect.
A service with a hidden multiplier violates that premise.

The reservation is important in two ways. It supplies a quantitative upper
bound on remaining authority, and it prevents replacement work from spending
the same capacity while the old effect remains possible. If a timeout could
refund $q_i$, another operation could consume it before the old one completes.
The safety account would then have forgotten part of its own exposure.
Conservative retention prevents that late-completion failure at the cost of
temporarily unavailable budget.

\subsection{Fenced revocation}
Fenced revocation first performs admission withdrawal. It then requests a
durable sink floor equal to the withdrawn issuer's new epoch. The sink serializes
the fence with acceptance and completion transactions. It cancels pending
permits below the floor and persists their terminal status. A completion already
inside its transaction either commits before the fence acquires the store, or
is cancelled by the fence before its own effect transaction. First deliveries
that arrive later receive durable cancellation tombstones.

The stronger acknowledgment is returned only after this sink transaction commits.
Let $a$ be its return. Every old permit requiring the withdrawn issuer has
request epoch strictly below the persisted floor. The verified fence predicate
then rejects it at completion. Any completion ordered before the fence has
already committed or cancelled before acknowledgment. Thus
\begin{equation}
 \text{old-authority effect charge strictly after }a=0.
\label{eq:v2fencedzero}
\end{equation}
This is an argument about the sink's transaction order. The core verifies the
epoch predicate and old-epoch exclusion lemma; the database ordering and effect
atomicity remain contract premises. The process tests exercise the ordering,
including a completion that holds the sink write lock while a fence waits.

The proof applies independently to every required issuer. An operator fence
blocks operator-required work; a steward fence blocks work requiring the
steward even if the operator remains live. A transfer with no steward requirement
does not become steward-controlled merely because that vector entry exists.
Renewal cannot lower the persisted floor or make a cleared old approval current.
Re-approval at a new epoch creates a distinguishable authority, rather than
resurrecting a withdrawn permit.

\subsection{Delay, availability, and ambiguous acknowledgment}
The guarantee after acknowledgment says nothing about a finite bound on when
acknowledgment becomes available. In the local service, a fence waits for an
in-progress SQLite writer. A remote service might be unavailable indefinitely.
If the fence response is lost, the caller does not infer success from elapsed
time; it retries or queries authoritative state while retaining uncertainty.
A monotone persisted floor makes repeated fence requests safe.

The experiments measure both acknowledgment delay and charge unavailable at
that boundary. In a delayed-completion control, the sink holds a transaction
for a declared 200 milliseconds; an independent connection confirms the write
lock before withdrawal and fencing. The fence waits behind that transaction.
The prior effect is allowed before acknowledgment, and subsequent retries add
no effects. The consumed charge remains unavailable because the effect
committed. In cancellation controls, terminal no-effect evidence permits release
after the stronger barrier. These are distinct availability outcomes.

This stronger boundary also explains why an ordinary process kill is
insufficient. Stopping the proposer closes its future computations but does not
establish any sink floor. Stopping the broker without recovery leaves its
accounting unavailable but does not cancel accepted work. Stopping the sink may
delay effects, but restarting it can resume durable pending requests. A
revocation acknowledgment must describe a durable authorization state change,
not merely a process lifecycle event.

\subsection{Unsupported contracts and limits}
The negative control applies the same operation twice at a service without an
idempotency identity. Both deliveries produce an effect. A correct reservation
ledger at the caller cannot retroactively make that service atomic or deduplicate
the deliveries. Nor can a client-only verifier guarantee a useful zero-effect
revocation acknowledgment if a noncooperating external service retains
unbounded latent work.

Multiple independent sinks would need a barrier covering every sink that could
still exercise the withdrawn authority. Partial acknowledgments could justify
partial exposure bounds, but not the global zero in
Equation~\ref{eq:v2fencedzero}. This artifact has one local cooperative sink and
does not evaluate distributed availability, partitions among several gateways,
or externally authenticated cancellation certificates. The mathematical
distinction between admission and fencing remains useful without claiming those
unimplemented extensions.
